\documentclass[11pt, a4paper]{scrartcl}
% Packages
\usepackage[margin=1.25in]{geometry}
\usepackage{index}
\usepackage{amsbsy} % Bold math symbols
\makeindex
\usepackage{tcolorbox}
\usepackage{varwidth}
\usepackage{amsmath, amssymb}
\usepackage{esint}
\usepackage{titlesec}
\usepackage{xcolor}
\usepackage{titling}
\usepackage[linktocpage]{hyperref}
\usepackage{pgfplots}
\pgfplotsset{compat=1.18}
\usepackage{caption}
\usepackage{amsthm}
\usepackage{import}
\usepackage{cancel}
\usepackage{caption}
\usepackage{nicematrix}
\usepackage{mathtools}
\usepackage{enumerate}
\usepackage{braket}
\usepackage{graphicx}
\usepackage{booktabs}
\usepackage{makecell}
\usepackage{tikz}
\usepackage[version=4]{mhchem}
\usetikzlibrary{arrows.meta, calc ,positioning, angles,quotes,patterns.meta,intersections}
\usepackage[italian]{babel}

% Font


% To reset footnote numbering each page
\usepackage[perpage]{footmisc}
\usepackage{setspace}
\setstretch{1.2}

% Titles 
\title{Fondamenti di\\ \vspace{.15cm} interazione radiazione-materia}
\author{}
\date{}


% svolgimento
\newenvironment{svolgimento}{\renewcommand\qedsymbol{$\blacksquare$}\begin{proof}[Svolgimento]}{\end{proof}}

\newtheoremstyle{style}% name of the style to be used
{5pt}% measure of space to leave above the theorem. E.g.: 3pt
{5pt}% measure of space to leave below the theorem. E.g.: 3pt
{\normalfont}% name of font to use in the body of the theorem
%{15pt}% measure of space to indent
{0pt}% measure of space to indent
{\noindent\bfseries}% name of head font
{}% punctuation between head and body
{ }% space after theorem head; " " = normal interword space
{\thmname{#1}\thmnumber{ #2}{\thmnote{ (#3)}.\ }}


\theoremstyle{style}
\newtheorem{esempio}{Esempio}[section]
\newtheorem{definizione}{Definizione}[section]
\newtheorem{prop}{Proposizione}[section]
\newtheorem{teorema}{Teorema}[section]
\newtheorem{lemma}{Lemma}[teorema]
\newtheorem{corollario}{Corollario}[teorema]
\newtheorem{osservazione}{Osservazione}[section]
\newtheorem{assunzione}{Assunzione}[section]
\newtheorem{notazione}{Notazione}[section]
\newtheorem{esercizio}{Esercizio}[section]

%% Generic box
\newtcolorbox{eqbox}[1][]
{
colback=gray!10,
arc=0pt,
boxrule=0pt,
title=#1
}

 \newenvironment{boxenv}[1][]{
    \begin{eqbox}[#1]
    }{
   \end{eqbox}
}

%Captions
\captionsetup[figure]{font=footnotesize,labelfont=footnotesize}
\captionsetup[table]{font=footnotesize,labelfont=footnotesize}
%Titlesec
\titleformat{\section}
{\fontsize{15}{20}\sffamily\scshape}
{\normalfont\color{gray}{\fontsize{20}{20}\selectfont\thesection}}
{0.7em}
{}
\hypersetup{colorlinks,breaklinks, linkcolor=[RGB]{74, 122, 164}}
\definecolor{asdf}{HTML}{4a7aa4}
% Personalizza la formattazione della subsection
\titleformat{\subsection}[block]{\fontsize{13}{20}\bfseries}{\normalfont\thesubsection}{.5em}{}


% Personalizza la formattazione della subsubsection
\titleformat{\subsubsection}[block]{\fontsize{12}{20}\bfseries}{\normalfont\thesubsubsection}{.5em}{}

% Maketitle customization
\renewcommand{\maketitle}{
\begin{center}
{\sffamily
{\fontsize{20}{20}\selectfont\MakeUppercase\thetitle}}

\vspace{0.2in}

{\large\scshape\sffamily\theauthor}
\end{center}
}


% Differenziale 
\newcommand{\dd}{\mathop{}\!\mathrm{d}}




%Evaluate symbol
\DeclareMathOperator{\di}{d\!}
\newcommand*\Eval[3]{\left.#1\right\rvert_{#2}^{#3}}

%%%%%%% Numero delle equazioni in formato a.b
\numberwithin{equation}{subsection}
%%%%%

%%%%%%%%%% Personalizzazione numeri lista
\renewcommand{\theenumi}{(\arabic{enumi})}

%%%% Table of contents

\usepackage[titles]{tocloft}

\renewcommand{\cftdot}{}
\usepackage{titletoc}
%\setcounter{tocdepth}{2}

%%%%%%%%%%%%%%%% Toc style

% Personalizzazione scritta indice


%%%% Lambda bar
\newcommand{\lbar}{{\mkern0.75mu\mathchar '26\mkern -9.75mu\lambda}}

\renewcommand{\arraystretch}{1.3}
\begin{document}
\maketitle
\newpage
\tableofcontents 
\newpage
\section{Ottica quantistica}

\subsection{Modi del campo elettromagnetico}

Nel seguito, si considereranno radiazioni elettromagnetiche solamente nel vuoto per parlare della quantizzazione del campo.
Dalle equazioni di Maxwell, allora, risulta
\begin{equation}
	\nabla ^2 \vec{E}(\vec{r},t) = \frac{1}{c^2} \frac{\partial ^2 \vec{E}(\vec{r},t)}{\partial t^2}  \hspace{2cm} \nabla \cdot \vec{E}(\vec{r},t) = 0
\end{equation}
Per trovare i modi del campo elettromagnetico, si risolvono queste due equazioni con condizioni al contorno date da una scatola cubica di lato $L$ per cui si richiedono soluzioni stazionarie sul bordo della scatola (equivalente a paretti conduttrici perfette).
Nel caso della componente lungo $\hat{x}$, per esempio, si trova un'espressione della forma 
\begin{equation}
	E_x(\vec{r},t) = E_x(t) \cos (k_x x ) \sin (k_y y ) \sin(k_z z)
\end{equation}
in modo che le componenti tangenziali si annullino.
Dalla condizione di annullamento al bordo, si ottiene che 
\begin{equation}
	k_x = \frac{\pi n_x}{L} \hspace{1cm} k_y = \frac{\pi n_y}{L} \hspace{1cm} k_z = \frac{\pi n_z}{L}
\end{equation}
dove $n_x,n_y,n_z \in [0,+\infty)$.
Nello spazio $(k_x,k_y,k_z)$, il numero di modi contenuti tra $k$ e $k +\dd k$ \`e dato dal volume infinitesimo $\dd V_k = 4 \pi k^2 \dd k$, diviso per la spaziatura tra un punto e l'altro, pari a $\pi / L$ per ciascuna direzione, che porta un fattore $(\pi / L ) ^{-3}$; avendo considerato solamente numeri d'onda positivi, si sta considerando una frazione pari a $\frac{1}{2^3}$ dello spazio originale e, in pi\`u, la presenza di due polarizzazioni per ciascun modo comporta un fattore $2$\footnote{Per $\vec{R}(\vec{r},t) = \vec{E}_0 e^{i (\vec{k}\cdot \vec{r}-\omega t)}$ e $\vec{k} = k \hat{z}$, per esempio, la condizione $\nabla \cdot \vec{E} = 0 \Rightarrow  \vec{k} \cdot \vec{E}_0 = 0$,  cio\`e $E_{0z}=0$. Questo vuol dire che $\vec{E}_0 = E_{0x}\hat{x} + E_{0y}\hat{y}$, cio\`e che per ogni vettore d'onda individuato da $k_z$, si hanno due modi indipendenti: uno lungo $\hat{x}$, ossia $E_{0x}$, e uno lungo $\hat{y}$, cio\`e $E_{0y}$.}.
Mettendo tutto insieme, la densit\`a di modi per unit\`a di volume \`e data dividendo il termine costruito sopra per $L^3$, cio\`e:
\begin{equation*}
	\rho (k)\dd k= \frac{ k^2 }{\pi^2}\dd k
\end{equation*}
Visto che $\omega = ck$, si ottiene 
\begin{equation}
	\rho (\omega ) \dd \omega = \frac{\omega ^2 \dd \omega}{\pi ^2 c^3}
\end{equation}
In questo modo, per sommare su tutti gli stati, invece di calcolare le somme sui $\vec{k}$ e sulle due polarizzazioni, si calcola $\int \frac{\omega^2}{ \pi^2 c^3} \dd \omega$, eventualmente aggiungendo il fattore $V$ per il volume.
\begin{osservazione}
Il considerare una densit\`a di modi per unit\`a di volume \`e utile per il fatto che, volendo studiare in seguito onde elettromagnetiche che non sono confinate in una regione di spazio, il termine $\rho (\omega ) \dd \omega$ \`e, cos\`i facendo, indipendente dal volume e, quindi, di validit\`a pi\`u generale.
\end{osservazione}

\noindent Sostituendo le soluzioni oscillanti accennate all'inizio, si osserva che la parte temporale soddisfa 
\[
\frac{\partial ^2 \vec{E}}{\partial t^2}  = - \omega ^ 2 \vec{E}(t) \implies \vec{E}(t) = \vec{E}(0) e^{-i\omega t}
\] 
Inoltre, l'energia del campo elettromagnetico \`e data, in MKS, da 
\begin{equation}
	\mathcal{E}  = \frac{1}{2} \int _V \Big[\varepsilon _0 \vec{E}^2(\vec{r},t) + \mu _0 ^{-1} \vec{B}^2 (\vec{r},t)\Big] \dd V
\end{equation}
dove si ricorda che $c = (\varepsilon _0 \mu _0)^{-1 / 2}$ e che
\[
	\nabla \times \vec{E}(\vec{r},t) = - \frac{\partial \vec{B}(\vec{r},t)}{\partial t} 
\] 
Si conclude che i campi soddisfano equazioni di oscillatori armonici con frequenza $\omega $ e hanno energia della stessa forma di quella di un oscillatore armonico, cio\`e \`e della forma $A p^2 + Bx^2$; infatti, si potrebbe pensare di associare al campo elettrico il ruolo di $x^2$, mentre al campo magnetico quello di $p^2$, o viceversa.
\begin{osservazione}
Dal discorso appena fatto, \`e ragionevole aspettarsi che, passando al caso quantistico introducendo degli operatori di campo per quello elettrico e quello magnetico, l'hamiltoniano del sistema abbia autoenergie che sono della stessa forma di quelle dell'oscillatore armonico, a meno del fatto che, invece di avere operatori impulso e posizione, si hanno i due operatori di campo.
\end{osservazione}

\noindent Dando per vera l'analogia nella precedente osservazione, si pu\`o vedere ogni singolo modo del campo elettromagnetico come un oscillatore armonico di frequenza $ \omega$ e un livello di eccitazione individuato da 
\[
E_n = \left(n + \frac{1}{2}\right) \hbar \omega \, , \ n = 0,1 , \ldots
\] 
In quest'ottica, $n$ rappresenta il numero di fotoni nel campo elettromagnetico: un fotone di frequenza $\omega$, quindi, \`e il quanto di eccitazione del modo del campo elettromagnetico relativo a tale frequenza.
\begin{osservazione}
Si nota la presenza di un'energia di vuoto $E_0 = \hbar \omega / 2$; questa non si misura negli esperimenti, in quanto si misura tipicamente l'eccitazione.
\end{osservazione}
\noindent La probabilit\`a (all'equilibrio termico) di trovare l'oscillatore armonico nel modo $n$-esimo, che secondo l'analogia corrisponde alla probabilit\`a di trovare $n$ fotoni nel campo elettromagnetico ($n$ fotoni in quel modo), \`e data da
\begin{equation}
	P(n) = \frac{1}{\mathcal{Z} } \exp \left(- \frac{E_n}{\kappa _B T}\right) 
\end{equation}
con $\mathcal{Z} = \sum_{n}^{} \exp\left(- \frac{E_n}{\kappa _B T}\right) $ funzione di partizione.
\begin{osservazione}
Pi\`u precisamente, questa corrisponde alla probabilit\`a di trovare un singolo oscillatore armonico eccitato al suo livello $n$-esimo, cio\`e di trovare $n$ fotoni nel suo modo.
\end{osservazione}
\noindent Chiamando per semplicit\`a $u = e^{- E_n / \kappa _B T}$ e usando 
\[
\sum_{n=0}^{+\infty} u ^ n = \frac{1}{1 - u}
\] 
allora $P(n) = (1-u) u ^n$.
Per questa ragione, il numero medio di fotoni nel modo $\hbar \omega$ \`e descritto da una Bose-Einstein con $\mu  =0 $:
\begin{equation}
	\begin{split}
		\langle N \rangle &= \sum_{n}^{} n P(n) = (1-u ) \sum_{n=0}^{+\infty} n U^n = (1-u) u \frac{\partial }{\partial u} \sum_{n=0}^{+\infty} u^n = \frac{u}{1-u}\\
				  & = \frac{1}{\exp \left(\frac{\hbar \omega}{\kappa _B T}\right) - 1}
	\end{split}
\end{equation}
Nel seguito, il numero medio di fotoni si indicher\`a con $\langle n \rangle$.
\subsubsection{Distribuzione di corpo nero}
La densit\`a di energia (termica) media \`e descritta dalla distribuzione di Plank per il corpo nero:
\begin{equation}
	\langle w_T(\omega ) \rangle \dd \omega = \langle n \rangle \hbar \omega \rho (\omega ) \dd \omega = \frac{\hbar \omega ^3}{\pi c^3} \frac{\dd \omega }{e^{\hbar \omega  / \kappa _B T }-1}
\end{equation}
L'espressione \`e data dal fatto che $\langle n \rangle\hbar \omega $ \`e l'energia totale media di un modo di frequenza $\omega$, cio\`e \`e il numero di fotoni che popolano tale modo per l'energia associata ad ogni quanto di eccitazione, il tutto moltiplicato per la densit\`a volumica di stati di frequenze tra $\omega$ e $\omega + \dd \omega $.
Il massimo di questa curva \`e raggiunto per $\omega _\text{max} \approx 2.8 \frac{\kappa _B T }{\hbar }$ e ha i seguenti andamenti asintotici:
\begin{align}
	&\langle w_T \rangle \sim \omega ^2 \tag{$\hbar \omega  \ll \kappa _B T $}\\
	& \langle w_T \rangle \sim \exp \left(- \frac{\hbar \omega }{\kappa _B T}\right)  \tag{$\hbar \omega \gg \kappa _B T$}
\end{align}
Integrando sulle frequenze, si trova la legge di Stefan-Boltzmann:
\begin{equation}
	\int _0^{+\infty} \langle w_T(\omega ) \rangle \dd \omega = \frac{\pi^2 (\kappa _B T)^4}{15 c^3 \hbar ^3} = \sigma T^4
\end{equation}

\subsubsection{Fluttuazioni}
Scrivendo $u$ in funzione di $n$, si trova che 
\[
u =  \frac{\langle n \rangle}{1+\langle n \rangle}
\] 
da cui 
\begin{equation}
	P(n) = \frac{\langle n \rangle ^n}{\big( 1 + \langle n \rangle\big) ^{1 + n}}
\end{equation}
Si vedr\`a che la probabilit\`a di avere $n$ fotoni in funzione del numero medio di fotoni nel modo ha questa stessa forma anche nel caso in cui il sistema non sia all'equilibrio termodinamico; in stato di non-equilibrio, cambier\`a la distribuzione di $\langle n \rangle$, ma l'espressione di $P(n)$ \`e la stessa.
La ``curva'' che ne risulta si vede essere sempre decrescente in $n$, quindi con $P(n=0)$ come il valore che \`e sempre pi\`u probabile, a prescindere da $\langle n \rangle$.
Si ha, inoltre, che decresce sempre meno rapidamente all'aumentare di $\langle n \rangle$.

La varianza di questa distribuzione si ottiene, al solito, come
\[
\sigma _n^2  = \sum_{n}^{} \left(n - \langle n \rangle\right) ^2 P(n) = \langle n^2 \rangle - \langle n \rangle^2
\] 
La distribuzione $P(n)$ \`e detta \textit{distribuzione geometrica} e si definiscono, in generale, i \textit{momenti fattoriali di ordine $r$} come $\langle n ( n-1) \ldots (n-r+1) \rangle$; nel caso delle distribuzioni geometriche, si sa che 
\[
\langle n ( n-1) \ldots (n-r+1) \rangle = r! \langle n \rangle^r
\] 
Si pu\`o usare questa relazione per concludere che $\langle n^2 \rangle - \langle n \rangle= \langle n (n-1) \rangle= 2 \langle n \rangle^2$, perci\`o
\begin{equation}
\sigma ^2_n = \langle n \rangle^2 + \langle n \rangle
\end{equation}
dove $\langle n \rangle^2$ si riconosce essere il contributo ``onda'', mentre $\langle n \rangle$ \`e il contributo ``particella''.
Di fatto, la varianza segue la parte ondulatoria quando $\langle n \rangle\gg 1$, cio\`e quando mediamente ci sono molte particelle, mentre segue il contributo particellare quando $\langle n \rangle \ll 1$, cio\`e quando la popolazione \`e molto piccola.

\subsection{Interazioni della radiazione con sistemi a due livelli}
Si considera un sistema a due livelli, di energie $E_1< E_2$ e popolazioni $N_1, N_2$, con $N = N_1+N_2$.
Si tratta un'interazione generica del sistema con il campo elettromagnetico, senza che questi siano all'equilibrio; questa \`e una situazione sperimentale tipica. 
In questo caso, la densit\`a di energia spettrale media si scompone in una parte termica e in una parte esterna: $\langle w(\omega ) \rangle = \langle w_T (\omega ) \rangle+ \langle w_E(\omega ) \rangle$.
Per il seguito, si assume che $\langle w_E(\omega ) \rangle$ vari molto lentamente con $\omega $ attorno alla frequenza di transizione $\frac{E_2 -E_1}{\hbar }$; questa assunzione \`e, in realt\`a, sempre verificata per $\langle w_T(\omega ) \rangle$.

La trattazione di Einstein prevede di considerare tre possibili comportamenti del sistema: emissione spontanea, indicata dal coefficiente $A_{21}$, emissione stimolata, indicata dal coefficiente $B_{21}$, e assorbimento, indicato a $B_{12}$. 
Visto che emissione stimolata e assorbimento coinvolgono l'interazione del sistema con un campo elettromagnetico e che l'intensit\`a della radiazione deve aumentare la probabilit\`a di transizione, si moltiplicano i coefficienti $B_{12}$ e $B_{21}$ per la densit\`a spettrale media di energia $w(\omega )$.
In questo modo, si hanno le \textbf{equazioni di rate}:
\begin{equation}
	\frac{d N_1}{d t}  = - \frac{d N_2}{d t}  = N_2 A_{21} - N_1 B_{12} \langle w(\omega ) \rangle + N_2 B_{21} \langle w(\omega ) \rangle
\end{equation}
Nel caso stazionario, in cui vale $\dot{N}_1 = \dot{N}_2 = 0$, e all'equilibrio termodinamico, quindi $\langle w_E(\omega ) \rangle = 0$, si ottiene
\begin{equation}
	\langle w_T(\omega ) \rangle = \frac{A_{21}}{\frac{N_1}{B_2}B_{12} - B_{21}}
\end{equation}
Essendo all'equilibrio termodinamico, si conosce il rapporto tra le popolazioni, visto che seguono la distribuzione di Boltzmann: assumendo che il livello $1$ abbia una degenerazione $g_1$ e che il livello $2$ abbia degenerazione $g_2$, si ha
\[
	\frac{N_1}{N_2} = \frac{g_1}{g_2} \overbracket{\exp \left(- \frac{E_1 - E_2}{\kappa _B T}\right) }^{= \exp \left(\frac{\hbar \omega}{\kappa _BT}\right) } \implies \langle w_T(\omega ) \rangle = \frac{A_{21}}{\displaystyle \frac{g_1}{g_2}\exp\left(\frac{\hbar \omega }{\kappa _B T}\right) B_{12}-B_{21}}
\] 
Per\`o questa deve avere la stessa forma della distribuzione di corpo nero, per cui devono valere le seguenti relazioni per i coefficienti:
\begin{equation}
	\begin{cases}
		g_1 B_{12} =g_2 B_{21}\\
		\\
		\displaystyle A_{21} = \frac{\hbar \omega ^3}{\pi^2 c^3}B_{21}
	\end{cases}
\end{equation}
Sostituendo le relazioni appena trovate in $\langle w_T(\omega ) \rangle$, si trova che
\begin{equation}
	\langle w_T(\omega ) \rangle = \frac{A_{21}\langle n \rangle}{B_{21}}
\end{equation}
In questo modo, volendo scrivere il rate di emissione totale per unit\`a di tempo, si trova che questo \`e dato da $B_{21} \langle w_T(\omega ) \rangle + A_{21} = A_{21}\left(\langle n \rangle + 1\right) $, dove l'uguaglianza \`e giustificata da quanto appena ricavato.
Ne segue che:
\begin{itemize}
	\item  se $\langle n \rangle\gg 1$, allora prevale l'emissione stimolata;
	\item se $\langle n \rangle \ll 1$, allora prevale l'emissione spontanea.
\end{itemize}
Visto che $\langle n \rangle$ \`e descritto dalla Bose-Einstein, si vede che $\langle n \rangle \sim 1$ quando $\frac{\hbar \omega }{\kappa _B T} \sim 0.7$. 
L'emissione stimolata sar\`a, allora, maggiormente rilevante a bassa frequenza, mentre a pi\`u alta frequenza prevarr\`a l'emissione spontanea.
\begin{osservazione}
A temperatura ambiente, $\kappa _B T \approx 25$ meV, per cui $\langle n \rangle \sim 1$ per $\hbar \omega \sim 10$ meV, corrispondenti ad un lontano infrarosso. 
\end{osservazione}

\subsubsection{Saturazione e attenuazione}

Si definisce \textbf{densit\`a di energia di saturazione} il coefficiente $w_s = \frac{\hbar \omega ^3}{\pi^2 c^3}$ in modo che $w_s B_{21} = A_{21}$; si osserva, per\`o, che questa \`e una densit\`a spettrale di energia.
La densit\`a di energia di saturazione, allora, \`e data da $w_s \dd \omega $.
	Tipicamente, per\`o, negli esperimenti non si usa questa quantit\`a, ma si usa l'intensit\`a (flusso di energia), che si trova essere legata alla densit\`a di energia tramite la relazione $I = c w$; si definisce, di conseguenza, l'\textbf{intensit\`a di saturazione} come $I_S = c w_s$, che, per\`o, \`e un'intensit\`a per unit\`a di frequenza.
	Integrando rispetto alle frequenze, si ottiene effettivamente l'intensit\`a totale relativa alla densit\`a di energia di saturazione.

\begin{esempio}
Si prende, come valore d'esempio, $\omega \sim 3 \cdot 10^{15}$ Hz.
Per queste frequenze, risulta una densit\`a di energia di saturazione dell'ordine di $w_s \dd \omega  \sim 10^{-14} \dd \omega \text{ J}/\text{m}^3$.
L'intensit\`a, invece, \`e dell'ordine di $I_S \dd \omega  \sim 3 \cdot 10^{-6} \dd \omega  \text{ W}/\text{m}^{-2}$.
Tenendo conto che, sperimentalmente, si hanno larghezze di riga dell'ordine di $10 ^{10}$ Hz, si ottiene che
	\[
	\int I_S \dd \omega \sim 10^{5}\, \frac{\text{W}}{\text{m}^{-2}}
	\] 
	In questo calcolo, si \`e usata l'assunzione fatta a inizio sezione, in cui si chiede che le densit\`a di energia varino, con $\omega $, molto lentamente attorno alla frequenza di transizione.
	In questo modo, l'integrale si pu\`o scrivere come $I_S(\omega _0) \Delta \omega $, dove $\Delta \omega $ \`e la larghezza di riga.
\end{esempio}

\begin{osservazione}
La differenza principale tra emissione spontanea e stimolata \`e che i primi sono emessi in modi e direzioni diverse, mentre i fotoni emessi per emissione stimolata mantengono le stesse identiche caratteristiche, direzione inclusa, del fotone che li ha fatti produrre.
\end{osservazione}

\noindent Si intuisce che l'emissione spontanea gioca un ruolo fondamentale nell'attenuazione di un fascio radiativo; per comprenderlo, si considera di irradiare un gas di atomi che non possono decadere dallo stato eccitato, se non per emissione stimolata e nel caso semplificato in cui $g_1=g_2\Rightarrow B_{12}=B_{21}=B$.
Assumendo che, nell'istante iniziale, tutti gli atomi siano nel fondamentale, questi inizieranno ad eccitarsi; tuttavia, il processo di assorbimento non smette di essere rilevante finch\'e le due popolazioni non si equivalgono:
\[
\dot{N}_1 = N_2 B \langle w(\omega ) \rangle - N_1 B \langle w(\omega ) \rangle \propto N_2 - N_1
\] 
Questo vuol dire che il numero di atomi eccitati continua a crescere finch\'e $\dot{N}_1 = 0$, che corrisponde alla condizione di stazionarit\`a: arrivati a questo punto, i due processi si equivalgono e viene emessa tanta radiazione, quanta ne viene assorbita (mediamente).
Le implicazioni di questo fenomeno stanno nel fatto che, per quanto gi\`a accennato, i fotoni emessi per emissione stimolata mantengono direzione e modo di quelli incidenti, quindi, in condizioni stazionarie, il materiale risulta trasparente.
In questo modo, se non sono presenti altri meccanismi di diseccitazione se non per emissione spontanea, non vi \`e alcuna attenuazione del fascio incidente.
Introducendo, invece, l'emissione spontanea, l'equazione in condizioni stazionarie diventa 
\[
\dot{N}_1 = 0 = N_2 A + (N_2 -N_1)B\langle w(\omega ) \rangle \implies \frac{N_2}{N_1} = \frac{B\langle w(\omega ) \rangle}{A + B \langle w(\omega ) \rangle} < 1
\] 
da cui $N_1 > N_2$ e, cio\`e, l'assorbimento prevale sull'emissione stimolata.
La frazione di energia persa per assorbimento, invece, si pu\`o ottenere come il numero di atomi che emettono per emissione spontanea, cio\`e $A_{21} N_2$, diviso il numero di atomi che emettono in totale, cio\`e $N_2 A_{21} + N_2 B_{21} \langle w \rangle$, che \`e pari a 
\[
\frac{A_{21}}{A_{21}+B_{21} \langle w \rangle} = \frac{1}{1 + \langle w \rangle / w_s}
\] 
\begin{osservazione}
In questo calcolo, si sta trascurando la piccola frazione di atomi che emettono spontaneamente nella direzione di propagazione del fascio.
\end{osservazione}

\noindent Dal risultato ottenuto, appare anche chiaro perch\'e si parli di densit\`a di energia di saturazione: all'incrementare di $\langle w \rangle$, l'attenuazione diventa un fenomeno trascurabile.

\subsubsection{Evoluzione delle popolazioni}

Per il seguito, si assume per semplicit\`a un sistema di due livelli, in cui quindi gli altri livelli siano molto distanti rispetto ai primi due, e si assume che $g_1 = g_2$ in modo che $B_{12} = B_{21} = B$ e si possano rimuovere gli indici.
Inoltre, si assume una $\langle w \rangle$ che sia uniforme all'interno della cavit\`a, cio\`e non vi sia attenuazione tale da rendere la distribuzione non-omogenea\footnote{Questa \`e una buona approssimazione nel caso in cui si invii fascio in una cavit\`a e la radiazione rimanga vincolata al suo interno, senza disperdersi.}.
Le equazioni di rate si scrivono come
\[
\frac{d N_1}{d t} = - \frac{d N_2}{d t}  = N_2 A + (N_2 - N_1) B \langle w \rangle
\] 
Il punto di stazionariet\`a si trova per $\dot{N}_1 = \dot{N}_2 = 0$, cio\`e 
\[
N_2 A + (N_2 - N_1) B \langle w \rangle = N_2 w_s + (N_2-N_1) \langle w \rangle = 0 
\] 
dove si \`e utilizzato che $A = w_s B$. 
\begin{center}
	\begin{minipage}
		{.45\columnwidth}
Da questa relazione, usando che $N_1 + N_2 = N$, si hanno le evoluzioni delle due popolazioni al variare di $\langle w \rangle$:
\[
\begin{cases}
	\displaystyle N_1 = \frac{w_s  + \langle w \rangle}{w_s + 2 \langle w \rangle}N\\
	\\
	\displaystyle N_2 = \frac{\langle w \rangle}{w_s + 2 \langle w \rangle}N\\
\end{cases}
\] 
	\end{minipage}
	\hfill
	\begin{minipage}
		{.5\columnwidth}
	\centering
	
\begin{tikzpicture}[y=0.039cm, x=.57cm]

    \def\N{100}
    \def\ws{1}

    % Assi
    \draw[->] (0,0) -- (10.5,0)
        node[right] {\footnotesize$\langle w\rangle$};
    \draw[->] (0,0) -- (0,110);

    % Riferimenti
    \draw (-0.1,\N) -- (0.1,\N)
        node[left] {$N$};
    \draw[dashed,gray] (0,50) -- (10,50);
    \node[left] at (0,50) {\footnotesize$N/2$};

    % Curva N_1
    \draw[cyan!60!black,thick]
        plot[domain=0:10,samples=200]
        (\x,{\N*(\ws+\x)/(\ws+2*\x)});

    % Curva N_2
    \draw[red!60!black,thick]
        plot[domain=0:10,samples=200]
        (\x,{\N*\x/(\ws+2*\x)});

        \node[cyan!60!black] at (5,65) {\footnotesize$N_1$};
        \node[red!60!black] at (5,35) {\footnotesize$N_2$};

\end{tikzpicture}
	\end{minipage}
\end{center}
Si sta considerando di avere una radiazione nel visibile, in modo da trascurare i fotoni termici, e che, in assenza di radiazione, tutti gli atomi siano nel fondamentale.
Si vede, allora, che all'aumentare della densit\`a di energia della radiazione, le due popolazioni tendono a diventare uguali, con una rapidit\`a che dipende da $w_s$.
Il massimo che si pu\`o fare, quindi, \`e rendere uguali le due popolazioni, ma non invertirle.

\begin{center}
	\begin{minipage}
		{.3\columnwidth}
\`E possibile, poi, determinare i rate di assorbimento emissione spontanea ed emissione stimolata, per unit\`a di $NA$, in funzione di $\langle w \rangle/w_s$; si nota che, asintoticamente, assorbimento e emissione stimolata sono rette.
\end{minipage}
	\hfill
	\begin{minipage}
		{.65\columnwidth}
	\centering
	\begin{tikzpicture}[
		scale=0.9,
    >=Stealth,
    thick,
    every node/.style={font=\small}
]

% Assi cartesiani
\draw[->] (0,0) -- (7.8,0)
    node[below right]
    {$ \langle\omega\rangle / \omega_s$};

\draw[->] (0,0) -- (0,4.2)
    node[above] {$R / NA$};

% Assorbimento
\draw[thick,domain=0:6.5,samples=150,smooth]
    plot (\x,{(\x*(1+\x))/(1+2*\x)})
    node[right] {Assorbimento};

% Emissione stimolata
\draw[thick,domain=0:6.5,samples=150,smooth]
    plot (\x,{(\x*\x)/(1+2*\x)})
    node[right] {Emissione stimolata};

% Emissione spontanea
\draw[thick,domain=0:6.5,samples=150,smooth]
    plot (\x,{\x/(1+2*\x)})
    node[right] {Emissione spontanea};

% Asintoto dell'emissione spontanea
\draw[dashed,thin] (0,0.5) -- (6.5,0.5);

\node[left] at (0,0.5) {$\frac12$};

% Differenza tra assorbimento ed emissione stimolata
\def\xpos{4.2}

\draw[<->,thin]
    (\xpos,{(\xpos*\xpos)/(1+2*\xpos)})
    --
    (\xpos,{(\xpos*(1+\xpos))/(1+2*\xpos)});
    \node at (4.5,1.7) {$R_{\mathrm{sp}}$};
    \draw[<->,thin](5,.5) -- (5,0);
    \node at (5.4,.2) {$R_{\mathrm{sp}}$};

% Origine
\node[below left] at (0,0) {$0$};

\end{tikzpicture}
	\end{minipage}
\end{center}
Si osserva, inoltre, che la distanza asintotica tra rate di emissione stimolata e assorbimento \`e pari al valore asintotico dell'emissione spontanea.

Si pu\`o valutare anche l'andamento temporale delle popolazioni, risolvendo l'equazione differenziale con condizioni iniziali $N_2(t=0) = 0$ e $N_1 ( t= 0) = N$; in questo modo, si ottiene
\[
	N_2(t) = N \frac{\langle w \rangle}{w_s + 2 \langle w \rangle} \bigg\{ 1 - \exp\Big[ - \big(A + 2 B \langle w \rangle\big) t \Big]  \bigg\} 
\] 
Quindi la popolazione $N_2$ raggiunge il valore di stazionariet\`a con una crescita esponenziale e una costante tempo data da $A + 2 B \langle w \rangle$.

Prendendo invece $N_2(t=0) = N$ e $N_1(t=0) = 0$ con $\langle w \rangle=0$, cio\`e in assenza di radiazione, allora si trova una decrescita esponenziale analoga degli atomi eccitati verso gli stati fondamentali per emissione spontanea.
\begin{osservazione}
Si nota che, nella trattazione, non solo la densit\`a di energia del campo elettromagnetico \`e media, ma anche i valori delle popolazioni $N_1$ e $N_2$.
Questi ultimi sono valori medi sia perch\'e varia il campo elettromagnetico, ma anche perch\'e si valutano tramite delle \textit{probabilit\`a di transizione}.
\end{osservazione}
\begin{boxenv}[]
\centering Riprendere dalla lezione 2 -10:30
\end{boxenv}
























\end{document}
